\section{Methods}
\label{sec:methods}
We address these questions through stakeholder interviews, transaction cost analysis, and a case study of the atmosfair-MiniGridCo PoAs in Nigeria.

\subsection{Research Design}
\label{subsec:design}
This study employs a qualitative research design using thematic analysis~\cite{BraunClarke} to investigate how voluntary carbon markets support rural solar mini-grid development in SSA. Compliance markets were excluded to maintain scope clarity. Qualitative approaches are particularly suited for exploring emerging carbon market applications where empirical data are scarce~\cite{CreswellPoth} and for generating theoretically grounded insights into structural barriers and enabling conditions.~\cite{Flyvbjerg}

Given limited empirical evidence on carbon-financed mini-grids in SSA (Section~\ref{subsec:problem}), the research adopts an exploratory approach focused on identifying conditions, barriers, and stakeholder experiences. The study examines carbon market integration across Kenya and Nigeria, selected to capture variation in regulatory frameworks, operational scales, and institutional arrangements. Rather than treating these as comparative cases, the analysis seeks patterns across contexts along three analytical dimensions: stakeholder type (developers, carbon experts, policy officials, registry professionals); carbon finance engagement level (pursuing, attempted, rejected); and project scale (predominantly 100--500~kWp).

This cross-cutting approach enables identification of themes recurring across multiple settings, including regulatory clarity, institutional capacity, aggregation mechanisms, and transaction cost economics, addressing two research questions: (RQ1) on commercial viability and (RQ2) on necessary conditions for integration at scale.

\textbf{Positionality:} The lead researcher brings professional experience in development finance and environmental assessment in SSA, providing contextual familiarity with regulatory frameworks and project economics. This proximity enhanced interpretive depth but required reflexive awareness of potential confirmation bias toward structural explanations, managed through an analytical journal maintained throughout the study.

\textbf{Limitations:} This qualitative design provides rich insights into conditions and barriers but does not establish causality or quantify impacts. Findings should be interpreted as preliminary identification of structural factors requiring validation through larger-scale quantitative research.

\subsection{Study Sites and Case Selection}
\label{subsec:site-selection}
Kenya and Nigeria were selected as illustrative contexts rather than comparative cases. Kenya has relatively advanced regulatory frameworks (e.g., the Energy Act of 2019 and the Carbon Markets Regulations of 2024) with 150 Kenya Off-Grid Solar Access Project (KOSAP)-deployed mini-grids but minimal carbon-financed activity, a gap that is itself analytically significant given its institutional development. Nigeria hosts approximately 200 mini-grids~\cite{rea2025} and includes the atmosfair-MiniGridCo PoA, one of the few operational carbon-financed frameworks in SSA, providing a critical case for examining aggregation feasibility.

Developer portfolios ranged from 16~kWp to 1~MWp, predominantly 100--500~kWp, reflecting the prevalence of small-to-medium scale mini-grids in SSA and enabling analysis of scale-dependent barriers such as transaction costs, methodological suitability, and revenue potential. 

\subsection{Data Collection}
This study received ethical approval from the SOAS University of London Research Ethics Panel, granted by the CeDEP Dissertation Convenor on 7 April 2025. Data collection combined semi-structured interviews and document analysis. Data were collected between June and August 2025.

\subsubsection{Semi-Structured Interviews}
\label{subsec:interviews}
Table~\ref{tab:participants} details the nine semi-structured interviews conducted with stakeholders spanning developers, carbon market experts, policy and regulatory officials, and registry professionals.

\begin{table}[t]
	\centering
	\resizebox{\textwidth}{!}{%
	\begin{tabular}{l c c c c}
		\toprule
		\textbf{Code} & \textbf{Stakeholder Type} & \textbf{Organization Type} & \textbf{Region} & \textbf{Duration} \\
		\midrule
		\textbf{DEV-KE-01} & Developer & Private sector & SSA & 45 min. \\
		\textbf{DEV-KE-02} & Developer & Private sector & SSA & 50 min. \\
		\textbf{DEV-NG-01} & Developer & Private sector & SSA & 35 min. \\
		\textbf{DEV-NG-02} & Developer & Private sector & SSA & 48 min. \\
		\textbf{EXP-01} & Carbon expert & Investor and Consultant & International & 60 min. \\
		\textbf{EXP-02} & Carbon expert & Carbon Marketplace & International & 38 min. \\
		\textbf{POL-01} & Regulatory policy expert & Regional professional & SSA & 45 min. \\
		\textbf{POL-02} & Regulatory policy expert & Regional professional & SSA & 45 min. \\
		\textbf{REG-01} & Registry expert & Gold Standard, ex-Verra & International & 75 min. \\
		\bottomrule
	\end{tabular}}%
	\caption{\textbf{Overview of Interview Participants.} All interviews conducted via virtual meetings.}
	\label{tab:participants}
\end{table}

Participants were contacted through multiple channels. Developers were identified through industry associations, including the Alliance for Rural Electrification, Gold Standard Registry, and snowball referrals.~\cite{bryman2016} Policymakers and carbon experts were approached through professional social media channels (e.g., LinkedIn) by analysing their experience relevance to the study. The carbon registry official was recruited via direct contact with the organization's technical team. Attempts to engage additional representatives from development finance institutions and consulting firms were unsuccessful within the research timeframe, limiting perspectives from investment and technical advisory communities.

\subsubsection{Document Analysis and Quantitative Data}
\label{subsec:doc-analysis}
Quantitative data were obtained from PoAs publicly available on the Gold Standard Registry, carbon certification cost data compiled from Gold Standard and Verra fee schedules and expert interviews, and project capital cost data obtained from developers and verified against the Africa Minigrid Developers Association (AMDA)~\cite{amda-bam} regional benchmarks. These secondary data provide contextual insight and should be considered indicative rather than precise measurements.

\subsection{Data Analysis}
\label{subsec:analysis}
Thematic analysis followed Braun and Clarke's~\cite{BraunClarke} six-phase framework, with coding conducted by the lead author consistent with reflexive thematic analysis principles.~\cite{BraunClarke2019} Data capture varied by participant consent: recorded interviews were transcribed verbatim, while non-recorded interviews were documented through structured notes capturing key themes and direct quotes where possible. All data were coded inductively in a structured Excel workbook tracking participant, extract, initial code, emerging theme, and research question alignment, enabling systematic cross-participant comparison and iterative refinement.~\cite{saldana}

Trustworthiness rested primarily on multi-source convergence: key themes were tested across nine interviews spanning four stakeholder types and documentary evidence. The transaction cost barrier, for example, emerged independently from developer interviews, was corroborated by expert assessment, confirmed by registry fee schedule analysis, and validated by Power Africa revenue leakage data.~\cite{power-africa-credits} An analytical journal tracked interpretive decisions throughout, and coding decisions were reviewed by the dissertation supervisor on a sample of transcripts. Negative cases were examined systematically. For instance, DEV-KE-01's projection of meaningful carbon revenues at scale qualified the transaction cost theme as scale-dependent rather than absolute.